% ===================================================================
%  TABEL Nr. 13 — Het hotel. — Het restaurant. — Het café.
%  Traduction 2026 d'apres texto/io, controlee sur texto/fr et
%  texto/es. Consigne : voir texto/nl/10-tabel-01.tex.
%
%  Deux notes, marquees toutes deux « (*) », et deux appels : celui du
%  bloc 01-2 (la chambre) et celui du bloc 09-1 (le menu).
% ===================================================================

%%K t13-noto-1 noto
\VUnotes{143.2mm}{%
\textit{\VUgras{(*) Zie voor de bijzonderheden van de kamer, tabel nr. 6.}}%
}

%%K t13-apar-1 apar
\VUsaut{3.0mm}
\VUtitre{15.0pt}{60}{TABEL Nr. 13}
\VUsaut{1.6mm}
\VUpk{10.2pt}{40}{Het hotel. --- Het restaurant.}
\VUsaut{2.0mm}
\VUpk{10.2pt}{40}{Het café.}
\VUsaut{3.4mm}

%%K t13-01-1 p
1. --- Nadat ik gisteravond in Royan aangekomen was, nam ik mijn intrek in het
« \textit{Hotel van de Oceaan} » dat mij door een vriend aanbevolen was.

%%K t13-01-2 p
Men gaf mij een mooie \VUgras{kamer} \textsuperscript{(2)}, op de tweede
\VUgras{verdieping} \textsuperscript{(1)} waar ik zeer goed sliep (*).

%%K t13-02-1 p
2. --- Vanmorgen, toen ik mijn kamer uitging, waarin de \VUgras{dienstbode}
\textsuperscript{(3)} het ontbijt gebracht had, ging ik de \VUgras{rookkamer}
\textsuperscript{(5)} binnen die ook als \VUgras{leeszaal}
\textsuperscript{(4)} gebruikt wordt om de \VUgras{kranten}
\textsuperscript{(6)} te lezen, terwijl ik een \VUgras{sigaret}
\textsuperscript{(8)} rookte. Na het lezen daalde ik af met de \VUgras{lift}
\textsuperscript{(9)} en ging door de stad wandelen.

%%K t13-03-1 p
3. --- Omdat ik vergeten was aan de \VUgras{bediende} \textsuperscript{(12)} van
het hotel te vragen op welk uur men de maaltijden aan de \VUgras{gemeenschappelijke
tafel} \textsuperscript{(11)} opdient, kwam ik vroeg terug en maakte ervan
gebruik om in de \VUgras{badkamer} \textsuperscript{(10)} van het hotel te gaan
baden. Daarna ging ik opnieuw naar de \VUgras{kassa} \textsuperscript{(15)} om
een van mijn vrienden op te bellen. Maar de \VUgras{telefoon}
\textsuperscript{(13)} was bezet; toen de \VUgras{caissière}
\textsuperscript{(14)}, die een \VUgras{rekening} \textsuperscript{(17)} in het
\VUgras{reizigersregister} \textsuperscript{(16)} inschreef, mijn verlegenheid
zag, vroeg zij de \VUgras{portier} \textsuperscript{(18)} de \VUgras{piccolo}
\textsuperscript{(19)} te sturen om mijn boodschap \VUgras{per fiets}
\textsuperscript{(20)} te doen.

%%K t13-04-1 p
4. --- Ik bedankte haar en ging naar buiten om een fooi te geven aan de
\VUgras{kruier} \textsuperscript{(24)} (de sjouwer) die uit het station op zijn
\VUgras{handkar} \textsuperscript{(25)} mijn \VUgras{hoedendoos}
\textsuperscript{(26)}, \VUgras{mijn valies} \textsuperscript{(27)} en
\VUgras{mijn reistas} \textsuperscript{(28)} gebracht had. De
\VUgras{brievenbesteller} \textsuperscript{(21)} die zojuist aangekomen was, gaf
mij een \VUgras{brief} \textsuperscript{(22)}.

%%K t13-05-1 p
5. --- Ik bleef nog enige tijd op de drempel van de deur, bij de
\VUgras{brievenbus} \textsuperscript{(23)}, om naar het verkeer op straat te
kijken. De \VUgras{conducteur} \textsuperscript{(30)} van de \VUgras{omnibus}
\textsuperscript{(29)} laadde op zijn rijtuig de \VUgras{koffer}
\textsuperscript{(31)} van een \VUgras{reiziger} \textsuperscript{(32)} van wie
de \VUgras{hotelhouder} \textsuperscript{(33)} afscheid nam. Een jonge
\VUgras{telegrambesteller} \textsuperscript{(35)} bracht zonder haast een
\VUgras{telegram} \textsuperscript{(34)} voor een gast van het hotel; een
\VUgras{toerist} \textsuperscript{(36)} met zijn \VUgras{fototoestel}
\textsuperscript{(38)} over de schouder, raadpleegde zijn \VUgras{reisgids}
\textsuperscript{(37)}.

%%K t13-tit-1 sub
\VUsaut{4.0mm}
\VUpk{10.2pt}{40}{Tijd voor het middagmaal.}
\VUsaut{3.0mm}

%%K t13-06-1 p
6. --- Voor het hek dommelde een zorgeloze \VUgras{boodschapper}
\textsuperscript{(39)} tussen zijn \VUgras{draaghaak} \textsuperscript{(40)} en
zijn \VUgras{poetsdoos} \textsuperscript{(41)}, terwijl een jonge
\VUgras{wasvrouw} \textsuperscript{(42)} die een \VUgras{mand}
\textsuperscript{(43)} met linnengoed droeg, lachte om de plechtige uitdrukking
van de \VUgras{gerant} \textsuperscript{(44)} die bij de deur stond.

%%K t13-07-1 p
7. --- Toen ik de \VUgras{kok} \textsuperscript{(45)} zag, die zijn
\VUgras{keuken} \textsuperscript{(47)} in het \VUgras{souterrain}
\textsuperscript{(48)} uitkwam om een \VUgras{vaatwasser}
\textsuperscript{(46)} een boodschap te laten doen, herinnerde ik mij dat het uur
van het middagmaal nabij is en ik ging het restaurant binnen, terwijl ik de
binnenplaats overstak waarop een \VUgras{automobilist} \textsuperscript{(50)}
zijn \VUgras{auto} \textsuperscript{(49)} stalde, terwijl een
\VUgras{monteur} \textsuperscript{(52)} volgens de aanwijzingen van de
\VUgras{fietser} \textsuperscript{(53)} een \VUgras{petroleumdriewieler}
\textsuperscript{(51)} herstelde die zojuist een ongeluk gehad had.

%%K t13-08-1 p
8. --- Ik vroeg aan een jonge \VUgras{piccolo} \textsuperscript{(54)} die
\VUgras{bierglazen} \textsuperscript{(55)} droeg, of de gemeenschappelijke tafel
onder de veranda staat, en op zijn bevestigend antwoord nam ik plaats aan een
van de tafels en liet mij een uitstekende maaltijd geven.

%%K t13-noto-2 noto
\VUnotes{143.2mm}{%
\textit{\VUgras{(*) Met behulp van de gerechtenlijst die in de woordenschat
opgenomen is, zal de onderwijzer het onderstaande menu gemakkelijk kunnen
variëren.}}%
}

%%K t13-tit-2 sub
\VUsaut{4.0mm}
\VUpk{10.2pt}{40}{Een uitstekend middagmaal.}
\VUsaut{3.0mm}

%%K t13-09-1 p
9. --- De kelner bracht mij het menu (*) van het middagmaal en de wijnkaart. Ik
koos en bestelde als \VUgras{voorgerecht garnalen} met \VUgras{boter}, en, als
tussengerecht, \VUgras{pens op de wijze van Caen}. Ik at geen \VUgras{vis}, maar
vroeg een portie \VUgras{lamsbout}, en, als \VUgras{groenteschotel, spinazie}
met room. Daarna, omdat geen van de \VUgras{tussengerechten} mij beviel, liet ik
verschillende nagerechten brengen, \VUgras{kaas}, \VUgras{vruchten} en
\VUgras{koekjes}. Ik voegde er bovendien een uitstekende
Saint-Émilion-\VUgras{wijn} aan toe, en zeer tevreden over mijn maaltijd
verliet ik de tafel nadat ik de \VUgras{kelner} \textsuperscript{(57)} een mooie
\VUgras{fooi} \textsuperscript{(59)} gegeven had.

%%K t13-10-1 p
10. --- Toen de \VUgras{beheerder} \textsuperscript{(60)} mij klaar zag om te
vertrekken, stelde hij mij voor naar het balkon te gaan om het prachtige
panorama van de \VUgras{Oceaan} \textsuperscript{(62)} te bewonderen. In de verte
begon men \VUgras{vissersbootjes} \textsuperscript{(63)}, \VUgras{zeilschepen}
\textsuperscript{(64)} en een enorme \VUgras{pakketboot} \textsuperscript{(65)}
te zien waarvan het zwarte silhouet zich tegen de witte \VUgras{wolken}
\textsuperscript{(67)} van de \VUgras{horizon} \textsuperscript{(66)} aftekent.
Een lichte bries deed de \VUgras{vlag} \textsuperscript{(70)} wapperen die boven
het \VUgras{uithangbord} \textsuperscript{(69)} van de \VUgras{hal}
\textsuperscript{(68)} bevestigd is.

%%K t13-tit-3 sub
\VUsaut{3.0mm}
\VUpk{10.2pt}{40}{Vermaak.}
\VUsaut{3.0mm}

%%K t13-11-1 p
11. --- Het schouwspel was zo boeiend, dat ik besloot morgen op het terras van
het café te klimmen met een \VUgras{toneelkijker} \textsuperscript{(71)} of een
\VUgras{verrekijker} \textsuperscript{(72)} mee.

%%K t13-12-1 p
12. --- Toen ik het balkon verliet, ging ik naar het café van het hotel om een
\VUgras{drankje} \textsuperscript{(73)} te drinken, terwijl ik een
\VUgras{partij biljart} \textsuperscript{(75)} speelde. Zonder stil te houden
stak ik de cafézaal over waarin veel gasten \VUgras{schaak}
\textsuperscript{(78)}, \VUgras{dammen} \textsuperscript{(79)},
\VUgras{domino} \textsuperscript{(80)}, \VUgras{triktrak}
\textsuperscript{(81)} of \VUgras{kaart} \textsuperscript{(82)} speelden, en ik
ging rechtstreeks naar de \VUgras{biljartzaal} \textsuperscript{(74)}.

%%K t13-13-1 p
13. --- Toen de partij afgelopen was, daalde ik naar de begane grond af en vroeg
de kelner een \VUgras{envelop} \textsuperscript{(84)}, \VUgras{papier}
\textsuperscript{(83)}, een \VUgras{postzegel} \textsuperscript{(85)}, een
\VUgras{pen} \textsuperscript{(87)} en \VUgras{inkt} \textsuperscript{(86)} om
een brief te schrijven.

%%K t13-13-2 p
Daarna riep ik een \VUgras{koetsier} \textsuperscript{(92)} wiens
\VUgras{rijtuig} \textsuperscript{(88)} voor het café was blijven staan. Ik vroeg
hem naar de prijs per uur of per rit, en, toen wij het over de prijs eens waren,
opende hij mij het \VUgras{portier} \textsuperscript{(89)} en installeerde mij.
Hij klom weer op de \VUgras{bok} \textsuperscript{(91)}, liet de paarden snel
vertrekken met een \VUgras{zweepje} \textsuperscript{(93)} en wij vertrokken voor
een bekoorlijke rit.
\VUsaut{6.0mm}
\VUfilet{22mm}
